\chapter{Viabilidad tecnica}

\section{Seguridad}

La viabilidad depende de que el sistema sea aceptable para seguridad industrial. Los AMR circularan con velocidades limitadas, zonas de proteccion configurables, escaner laser de seguridad, bumpers, parada de emergencia y protocolos de retirada manual. La validacion debe incluir pruebas con peatones, obstaculos inesperados y estaciones ocupadas.

\section{Integracion en planta existente}

La integracion puede realizarse por fases. En el piloto, las misiones se pueden cargar manualmente desde HMI y registrar en una base de datos local. En la expansion, el sistema debe intercambiar ordenes con un planificador o MES. En el despliegue extendido, se recomienda integracion bidireccional con sistemas de produccion, mantenimiento y calidad.

\section{Limitaciones y riesgos}

\begin{itemize}
  \item El AMR no elimina la necesidad de planificacion de materiales ni de disciplina operacional.
  \item La inspeccion FOD con camara es preventiva y complementaria; no sustituye procedimientos certificados.
  \item La red WiFi industrial y la calidad del mapa condicionan la disponibilidad.
  \item La flota requiere mantenimiento de baterias, ruedas, sensores y calibraciones.
  \item La aceptacion de los operarios puede caer si el sistema se percibe como obstaculo o carga administrativa.
\end{itemize}

\section{Mitigaciones}

Las mitigaciones incluyen piloto acotado, formacion practica, participacion temprana de operarios, rutas validadas en CoppeliaSim, pruebas de seguridad, mantenimiento preventivo y despliegue por etapas. Para no depender de una integracion compleja desde el primer dia, el sistema se diseña con una base de datos propia y conectores progresivos hacia MES/ERP.

\section{Escalabilidad}

La propuesta es escalable porque separa robot, gestor de flota, base de datos y supervision. El mismo modelo de mision puede aplicarse a un robot o a una flota, ajustando reglas de prioridad, zonas y ventanas de entrega. Esta arquitectura es coherente con la evolucion hacia sistemas multi-robot descrita en la literatura de robotica movil \cite{kagan2019}.
